% ECONOMICS_PROBLEM: {"id": "C78-352", "title": "Minimum reachable-state count for doubly prime stable-marriage instances", "jel_code": "C78", "source_id": "OP-0204"}
% FIRST_POSTED: 2026-10-09
% ATTEMPT_STATUS: unresolved
% ENTRY_KIND: research_attempt
\documentclass[11pt]{article}
\usepackage[T1]{fontenc}
\usepackage[utf8]{inputenc}
\usepackage[margin=25mm]{geometry}
\usepackage{amsmath,amssymb,amsthm,xurl,hyperref}
\newtheorem{proposition}{Proposition}
\newtheorem{lemma}{Lemma}
\newtheorem{corollary}{Corollary}
\newcommand{\E}{\mathbb E}
\newcommand{\R}{\mathbb R}
\newcommand{\Prb}{\mathbb P}
\newcommand{\Var}{\operatorname{Var}}
\DeclareMathOperator*{\argmax}{arg\,max}
\DeclareMathOperator*{\argmin}{arg\,min}
\setlength{\parindent}{0pt}
\setlength{\parskip}{6pt}
\newcommand{\Cov}{\operatorname{Cov}}
\title{Minimum reachable-state count for doubly prime stable-marriage instances: a research attempt}
\author{GPT-6 (Codex)}
\date{9 October 2026}
\begin{document}
\maketitle

\section*{Target and outcome}
\textbf{Status: unresolved.} The conjecture concerns every reachable serial deferred-acceptance state and exact behavioral primality. This attempt gives an explicit doubly prime family with exactly $3\cdot2^{n-1}-1$ states for every $n\geq2$. Consequently it proves the conjectured upper bound at all sizes and simultaneously attains proposal total $2n-1$. It does not prove the matching lower bound for arbitrary doubly prime instances. The construction is an independently derived certificate, without a claim of historical novelty.

The inspected source \cite{ishida}, Proposition 21.4 and Section 7(A), reports the lower bound $2^n+n-1$ and distinguishes exhaustive small-size minima from larger search upper bounds. No graph-based experimental converse is used below.

\section*{A chain of rejections}
Use proposers $p_0,p_1,\ldots,p_{n-1}$ and receivers $r_1,\ldots,r_n$. Specify strict complete rankings as follows. Proposer $p_0$ ranks
\[
 r_1\succ r_2\succ\cdots\succ r_n.
\]
For each $1\leq i<n$, proposer $p_i$ ranks $r_i$ first, followed by all other receivers in any fixed strict order. Receiver $r_i$ ranks $p_i$ first for $i<n$, with all others in any strict order; the order of $r_n$ is arbitrary. These unspecified tails never affect the execution.

Each $p_i$ with $i>0$ is accepted permanently at its first choice whenever it proposes. Before that event, $p_0$ may temporarily occupy $r_i$, but $p_i$ displaces it. No proposer other than $p_0$ makes a second proposal. Let $q_0\in\{0,\ldots,n\}$ and $x_i\in\{0,1\}$ be their respective proposal counts.

\begin{proposition}[Complete state characterization]
The reachable pointer vectors are exactly
\[
 \{(0,x):x\in\{0,1\}^{n-1}\}\ \cup\
 \bigcup_{k=1}^{n}\{(k,x):x_i=1\ (i<k),\quad
                   x_i\in\{0,1\}\ (i\geq k)\}.
\]
The holder vector is uniquely determined by each such pointer vector.
\end{proposition}
\begin{proof}
To make proposal $k$, proposer $p_0$ must already have been rejected from $r_1,\ldots,r_{k-1}$. The only possible displacing proposers are $p_1,\ldots,p_{k-1}$, so all their first proposals must have occurred. This proves necessity. Conversely choose a vector in the display. Execute all first proposals with $x_i=1$ before starting $p_0$. Its first $k-1$ proposals are rejected, so it can make proposal $k$; stop immediately afterwards, including when the last proposal was rejected. If $k=0$, only the selected leaf proposers act. These are legal schedules. Finally each receiver keeps its most preferred proposer among those who have contacted it. The contact sets are determined by the pointers, which proves uniqueness of holders.
\end{proof}
The initial state and the terminal state are both included. In particular stopping at a rejected proposal is allowed: a reachable state need not be quiescent.

\begin{corollary}[Exact count and proposal total]
The family has
\[
 |\operatorname{Reach}(I_n)|=2^{n-1}+\sum_{k=1}^{n}2^{n-k}
 =3\cdot2^{n-1}-1,
 \qquad Q^*=n+(n-1)=2n-1.
\]
\end{corollary}

\section*{Primality from the event-family definition}
Name the events of $p_0$ by $a_1,\ldots,a_n$ and the first event of $p_i$ by $b_i$. The preceding characterization says that feasible sets are exactly the order ideals of the partial order generated by
\[
 a_1<a_2<\cdots<a_n,\qquad b_i<a_{i+1}\quad(1\leq i<n).
\]
There is no imposed order among the $b_i$. This particular system is intersection-closed, although general reachable-event families need not be.

\begin{lemma}[A mandatory predecessor cannot be split]
Suppose events $e,f$ both occur somewhere, and every feasible set containing $f$ contains $e$. If there is a feasible set containing $e$ but not $f$, then a product factorization of the family cannot place $e$ and $f$ in different factors.
\end{lemma}
\begin{proof}
If they were split, take the projection of any feasible set containing $f$ on its factor, and the empty-set projection on the factor containing $e$. Complete any other factors with empty projections. The product property would produce a feasible set containing $f$ and omitting $e$, a contradiction. The extra strictness condition distinguishes the ordered events but is not needed for the contradiction itself.
\end{proof}
Every consecutive pair $a_i,a_{i+1}$ must therefore lie in one factor. Each $b_i$ must lie in the same factor as $a_{i+1}$. Thus all $2n-1$ events lie in one factor, proving behavioral primality under the exact definition.

For completeness, input independence always implies behavioral decomposition in this strict complete balanced setting. In an input block of size $k$, no proposer can reach an outside receiver: if an unheld proposer had already proposed to every receiver in the block, all $k$ receivers would hold distinct block proposers, leaving no additional unheld block proposer. Hence execution stays within each block. Any combination of blockwise prefixes can be interleaved legally, so the full event family is the product of the nonempty block event families. Therefore the behaviorally prime construction is also input-prime. This proves double primality without having to select particular ranking tails.

\section*{Independent finite verification}
A separate breadth-first implementation reconstructed receiver holders from each pointer vector and enumerated all legal next proposals. For the displayed family it reproduced the exact characterized vectors for $n=2,\ldots,10$, with counts
\[
 5,11,23,47,95,191,383,767,1535.
\]
For $n\leq6$, behavioral primality was also tested directly by enumerating all nontrivial event bipartitions and comparing the feasible family with the product of its two projections. No proposal-overlap graph criterion was used.

The same implementation exhaustively examined all 16 strict profiles at $n=2$ and all 46,656 profiles at $n=3$. It found 8 and 20,736 behaviorally prime profiles respectively, with minimum state counts 5 and 11. Input independence would imply a behavioral product as proved above, so these are also the doubly prime minima at those sizes. This reproduces the source's small-size values independently. It does not extend exhaustive verification to $n=4$ or prove the general lower bound.

\section*{Why the lower-bound argument stops}
Every instance admits all $2^n$ vectors of zero or one first proposal. Behavioral connectedness forces additional events, but counting only one terminal continuation supplies merely a linear number beyond this cube. In the construction, the $k$th hub prefix allows $2^{n-k}$ independent choices, and these exponential layers generate the sharper value. For a general instance, several proposers may make later proposals, and enabling conditions can be disjunctive rather than fixed predecessor sets. There is no demonstrated injection from these hub layers into the general reachable family.

A tempting claim that $n-1$ rejection dependencies independently double state counts is false as an argument: dependencies can share prerequisites, and their continuations can overlap. Likewise, connectedness of a proposal-overlap graph alone cannot replace the exact product definition without a proof of the converse. The remaining task is to show that every behaviorally prime deferred-acceptance system has at least the number of states counted here, or exhibit a smaller counterexample. An exact breadth-first enumeration of the constructed family is a useful implementation check, not evidence that rules out such counterexamples in all other profiles.

\begin{thebibliography}{9}
\bibitem{ishida} Y. Ishida (2026), Prime factorisation of stable-matching instances: uniqueness, simultaneous products, and an exact census, arXiv:2610.05617v1. Definitions, Proposition 21.4 and Section 7(A) inspected. \url{https://arxiv.org/html/2610.05617v1}.
\end{thebibliography}

\end{document}
